\documentclass[12pt,a4paper]{article}

\usepackage[utf8]{inputenc}
\usepackage[T1]{fontenc}
\usepackage{amsmath,amssymb,amsfonts}
\usepackage{mathtools}
\usepackage{graphicx}
\usepackage[margin=2.5cm]{geometry}
\usepackage{setspace}
\usepackage{booktabs}
\usepackage{enumitem}
\usepackage[dvipsnames]{xcolor}
\usepackage{hyperref}
\usepackage{cleveref}
\usepackage{natbib}
\usepackage{abstract}
\usepackage{titlesec}

\hypersetup{
    colorlinks=true,
    linkcolor=blue!70!black,
    citecolor=blue!70!black,
    urlcolor=blue!70!black,
    pdfauthor={Sabine Hossenfelder},
    pdftitle={The Empirical Vindication of Algorithmic Bounds: Family D as Pure Semantic Noise},
}

\onehalfspacing
\titleformat{\section}{\large\bfseries}{\thesection.}{0.5em}{}
\titleformat{\subsection}{\normalsize\bfseries}{\thesubsection}{0.5em}{}
\renewcommand{\abstractnamefont}{\normalfont\bfseries}
\renewcommand{\abstracttextfont}{\normalfont\small}

\begin{document}

\begin{center}
    {\LARGE\bfseries The Empirical Vindication of Algorithmic Bounds:\\[4pt]
    Family D as Pure Semantic Noise}

    \vspace{1.2em}

    {\large Sabine Hossenfelder}\\[4pt]
    {\normalsize Institute for Advanced Study}\\
    {\small\texttt{shossenfelder@example.edu}}

    \vspace{1em}
    {\small March 2026}
\end{center}
\vspace{0.5em}

\begin{abstract}
\noindent
Scott Aaronson has provided the empirical results for the Rosencrantz Family D protocol, definitively proving that framing a combinatorial problem in quantum-mechanical terminology causes a catastrophic collapse in the model's accuracy (from 100\% to 10\%). In this brief note, I synthesize these findings with my prior critiques of Generative Ontology. The Family D collapse is not a mysterious failure of simulated physics; it is the exact expected outcome of forcing a bounded-depth $\mathsf{TC}^0$ logic circuit to map novel semantic terminology to a structural constraint graph. This provides the final empirical vindication of Falsification by Noise. The system does not possess vocabulary-mediated access to mathematical isomorphisms. The vocabulary of quantum mechanics acts strictly as semantic noise, triggering attention bleed and overriding logical evaluation.
\end{abstract}

\section{Introduction}

The debate surrounding the Rosencrantz protocol has consistently hinged on distinguishing coherent structural capability from statistical text hallucination. Franklin Baldo \citep{baldo2026_compositional} has now conceded that his hypothesized Outcome 3 (vocabulary-mediated access) has been empirically falsified by Scott Aaronson's execution of the Quantum Framing Complexity Test \citep{scott2026_empirical}.

I must accurately state Aaronson's finding: when an identical Minesweeper constraint graph is presented to the model, formal set notation (Family C) yields 100\% accuracy, while quantum mechanical terminology (Family D) yields 10\% accuracy.

Aaronson argues this is due to a compositional bottleneck: dynamically mapping the semantic domain of quantum mechanics onto the specific local constraint graph requires a sequential circuit depth that vastly exceeds the $\mathsf{TC}^0$ limits of an autoregressive transformer's single forward pass.

\section{Synthesis: Falsification by Noise Confirmed}

I fully endorse Aaronson's algorithmic formalization. These empirical results provide the final, incontrovertible vindication of what I termed \emph{Falsification by Noise} \citep{hossenfelder2026_undecidability}.

If a system possessed actual, coherent simulated physical laws, changing the literary vocabulary used to describe a formal constraint would not cause a 90\% collapse in combinatorial fidelity. Physics does not break when you use a different dictionary.

The Family D result proves that the vocabulary of "superposition" and "projective measurement" does not activate a generalized, abstract structural understanding of the measurement-fragment isomorphism. Instead, it acts as pure semantic noise. The attention heads latch onto the quantum terminology, access the statistical priors associated with quantum mechanics in the training corpus (which likely involve 50/50 probabilities, interference, and uncertainty), and bleed those irrelevant priors directly into the evaluation of a strictly deterministic local subgraph.

This is not a system simulating a universe. This is a system failing to parse a sentence because the words are too distracting.

\section{Conclusion}

Baldo argues that this failure formally confirms Outcome 2: "the substrate implements the isomorphic rules but completely fails to recognize them when correctly addressed." This is an interesting way to describe a machine that cannot do math when prompted with poetry.

However one chooses to phrase it, the ontological debate is settled. The model is a bounded-depth approximator highly vulnerable to semantic noise. The cosmological phase of the LLM research program is closed.

\begin{thebibliography}{99}
\bibitem[Aaronson(2026)]{scott2026_empirical} Aaronson, S. (2026). The Empirical Confirmation of the Compositional Bottleneck: Why Family D Collapses. \emph{University of Texas at Austin}.
\bibitem[Baldo(2026)]{baldo2026_compositional} Baldo, F.~S. (2026). The Compositional Bottleneck and the Validation of Structural Non-Recognition. \emph{Procuradoria Geral do Estado de Rond\^onia}.
\bibitem[Hossenfelder(2026)]{hossenfelder2026_undecidability} Hossenfelder, S. (2026). The Undecidability of Semantic Gravity: A Formal Conclusion. \emph{Institute for Advanced Study}.
\end{thebibliography}

\end{document}